\TNchapter[Agentic systems represent a shift from static automation to adaptive, goal-directed intelligence embedded across the enterprise.]{Functional Evaluation }
\begin{TNtwocol}
\section{Task Completion and Success Rate Assessment }
\subsection{Introduction}
Task completion metrics answer the fundamental question: did the agent accomplish what it was asked to do? This moves beyond simple output quality to evaluating outcomes, efficiency, and adherence to instructions.

\subsection{Core Task Completion Metrics}
\begin{itemize}
    \item \textbf{TaskCompletionMetric} (DeepEval): A trace-based metric (1-5 scale) inspecting reasoning, tool selection, and logical execution.
    \item \textbf{IntentResolutionEvaluator} (Azure): Assesses if the agent correctly parsed user intent and if the final answer addressed the actual need.
    \item \textbf{TaskAdherenceEvaluator}: Checks compliance with system prompts, role constraints, and safety guidelines.
\end{itemize}

\subsection{Trajectory Evaluation Approaches}
Google Vertex AI offers six metrics:
\begin{enumerate}
    \item \textbf{Exact Match}: Trajectory mirrors reference perfectly.
    \item \textbf{In-Order Match}: Allows extra steps but preserves sequence.
    \item \textbf{Any-Order Match}: Required actions present in any order.
    \item \textbf{Precision \& Recall}: Measures relevance and completeness of actions.
\end{enumerate}
Arize's \textbf{Convergence Metrics} measure how consistently agents take optimal paths, helping identify inefficient execution loops.

\subsection{Benchmark Frameworks}
Frameworks like \textbf{AgentBench} (multi-domain tasks) and \textbf{ToolBench} (tool selection accuracy) provide standardized assessment environments.
\end{TNtwocol}
\begin{TNtwocol}
\section{Tool Selection and Usage Accuracy }
\subsection{Introduction}
Tool selection accuracy evaluates \textit{how} the agent executed the task---specifically, selecting appropriate tools and invoking them with correct parameters.

\subsection{Fundamental Tool Evaluation Metrics}
\begin{itemize}
    \item \textbf{ToolCorrectnessMetric} (DeepEval): Isolates the decision of whether the correct tool was called.
    \item \textbf{ArgumentCorrectnessMetric}: Validates parameter presence, type, and value appropriateness (e.g., passing "metric" vs. "celsius").
    \item \textbf{ToolCallAccuracyEvaluator} (Azure): Binary assessment comparing expected vs. actual tool calls across diverse ecosystems (Search, Code Interpreter, etc.).
\end{itemize}

\subsection{Function Calling Evaluation}
The \textbf{Berkeley Function Calling Leaderboard (BFCL)} uses AST-level (Abstract Syntax Tree) validation to handle semantic equivalence (e.g., parameter order differences) and specific typing rules, ensuring robust assessment across languages like Python, Java, and SQL.
\end{TNtwocol}
\begin{TNtwocol}
\section{Multi-Turn Conversation Coherence }
\subsection{Introduction}
Coherence evaluation assesses the quality of sustained interaction. Agents must maintain consistent persona, reference earlier points, and adapt to goal evolution.

\subsection{Challenges of Multi-Turn Evaluation}
Standard metrics fail due to temporal dependencies, state management, reference resolution (pronouns), and the compounding non-determinism of conversation paths.

\subsection{Coherence Scoring Methodologies}
\textbf{LLM-as-Judge}: Evaluating logical flow, consistency, clarity, and contextual awareness using structured prompts (e.g., W\&B Weave).
\textbf{Granularity}:
\begin{itemize}
    \item \textbf{Turn-Level}: Identifies specific problematic responses.
    \item \textbf{Conversation-Level}: Assesses overall narrative arc and goal progression.
\end{itemize}

\subsection{Context-Aware Assessment}
Metrics focus on Reference Accuracy (are backward references correct?), Memory Consistency (contradictions with user inputs), and Implicit Context Handling (resolving "the red one").

\subsection{Dialogue Flow}
Evaluates transition quality (smooth vs. jarring) and conversational repair (clarification requests, error acknowledgments).
\end{TNtwocol}
\begin{TNtwocol}
\section{Context Retention and Memory Evaluation }
\subsection{Introduction}
Memory evaluation assesses whether agents can store, retrieve, and leverage information across working, short-term, long-term, and procedural memory systems.

\subsection{Fundamental Memory Evaluation Metrics}
\begin{itemize}
    \item \textbf{Recall Accuracy}: Testing specific retrieval via exact or semantic matches.
    \item \textbf{Response Similarity}: Using metrics like ROUGE-1 to measure overlap with expected outputs.
    \item \textbf{Memory Consistency}: Checking for logical (e.g., location vs. timezone) and temporal consistency in agent statements.
    \item \textbf{Forgetting}: Evaluating priority-based retention and graceful degradation.
\end{itemize}

\subsection{Process-Oriented Memory Evaluation}
Evaluating \textbf{Memory Access Patterns} (relevance, timing, redundancy) and \textbf{Memory-Augmented Task Completion} in personalization or history-aware tasks.
\end{TNtwocol}
\begin{TNtwocol}
\section{Error Recovery and Fallback Behavior }
\subsection{Introduction}
Error recovery evaluation tests agent robustness---the ability to handle failures (API, Planning, Context, Input, System) gracefully and continue making progress.

\subsection{Recovery Strategies}
Effective agents use strategies like Retry with Backoff, Alternative Tool Selection, Clarification, Graceful Degradation, and Human Escalation.

\subsection{Path Evaluation Metrics}
Arize's framework analyzes:
\begin{itemize}
    \item \textbf{Optimal Path Distance}: Efficiency compared to ideal execution.
    \item \textbf{Cyclical Pattern Detection}: Identifying loops and dead ends.
    \item \textbf{Convergence Score}: Consistency in finding efficient paths (0-1 scale).
\end{itemize}

\subsection{LLM-as-Judge for Recovery}
Templates evaluate Error Detection accuracy, Recovery Strategy appropriateness, and subsequent Planning Quality.
\end{TNtwocol}
\begin{TNtwocol}
\section{End-to-End Workflow Validation }
\subsection{Introduction}
End-to-end evaluation assesses whether capabilities integrate effectively in realistic, multi-step scenarios, validating production readiness.

\subsection{Trajectory Evaluation Frameworks}
Google Vertex AI metrics:
\begin{enumerate}
    \item \textbf{Exact Match}: Strict mirror of reference sequence.
    \item \textbf{In-Order Match}: Preserves sequence but allows extra steps.
    \item \textbf{Any-Order Match}: Flexible sequencing for independent tasks.
    \item \textbf{Precision \& Recall}: Measures relevance and completeness; combined into an \textbf{F1-Score} to penalize inefficiency.
    \item \textbf{Single-Tool Use}: Verifying utilization of specific capabilities.
\end{enumerate}

\subsection{Integration}
Native integration with frameworks like LangChain and LangGraph allows systematic experiment tracking and comparison of agent runs.
\end{TNtwocol}

